\begin{problem}[2019年第36届全国中学生物理竞赛复赛][电磁轨道炮简化模型]{106}
某电磁轨道炮的简化模型如图a所示，两圆柱形固定导轨相互平行，其对称轴所在平面与水平面的夹角为 $\theta$，两导轨的长均为 $L$、半径均为 $b$、每单位长度的电阻均为 $\lambda$，两导轨之间的最近距离为 $d$（$d$ 很小）。一质量为 $m$（$m$ 较小）的金属弹丸（可视为薄片）置于两导轨之间，弹丸直径为 $d$、电阻为 $R$，与导轨保持良好接触。两导轨下端横截面共面，下端（通过两根与相应导轨同轴的、较长的硬导线）与一电流为 $I$ 的理想恒流源（恒流源内部的能量损耗可不计）相连。不考虑空气阻力和摩擦阻力，重力加速度大小为 $g$，真空磁导率为 $\mu_{0}$。考虑一弹丸自导轨下端从静止开始被磁场加速直至射出的过程。
\begin{subproblem}
求弹丸在加速过程中所受到的磁场作用力；
\end{subproblem}
\begin{subproblem}
求弹丸的出射速度；
\end{subproblem}
\begin{subproblem}
求在弹丸加速过程中任意时刻、以及弹丸出射时刻理想恒流源两端的电压；
\end{subproblem}
\begin{subproblem}
求在弹丸的整个加速过程中理想恒流源所做的功；
\end{subproblem}
\begin{subproblem}
在 $\theta = 0^{\circ}$ 的条件下，若导轨和弹丸的电阻均可忽略，求弹丸出射时的动能与理想恒流源所做的功之比。
\end{subproblem}
\end{problem}